%% ma: L10-B12-0066
%% lop: 10
%% chuong: 5
%% bai: 12
%% dang: 10.12.4
%% loai: DS
%% muc: [TH, NB, TH, TH]
%% dap_an: "ĐSĐĐ"
%% nguon: "chuyen-de-cac-so-dac-trung-cua-mau-so-lieu-khong-ghep-nhom-toan-10-chuong-trinh-moi.pdf (GV. Phan Nhật Linh) (P01), Bài 01 Phần II Câu 10"
%% kien_thuc: "Lớp 10 Bài 12 (số gần đúng, sai số tuyệt đối, độ chính xác, sai số tương đối, quy tròn số)"
%% trang_thai: chinh-thuc
%% ghi_chu: "Sửa lời các ý a), b), c) cho chặt (gốc viết 'có sai số tương đối là'). Tự duyệt theo yêu cầu giáo viên 22/10/2026."
\begin{ex}
Ba học sinh An, Ba, Na lần lượt đo chiều cao của mình. Bạn An đo được $168\text{ cm}\pm0{,}5\text{ cm}$, bạn Ba đo được $181\text{ cm}\pm2\text{ cm}$, bạn Na đo được $148\text{ cm}\pm1\text{ cm}$. Xét tính đúng sai của các khẳng định sau:
\choiceTF{\True Sai số tương đối của phép đo của bạn An không vượt quá $\dfrac{0{,}5}{168}$}{Sai số tương đối của phép đo của bạn Ba không vượt quá $\dfrac{-2}{181}$}{\True Số $\dfrac1{148}$ (chặn trên của sai số tương đối phép đo của bạn Na) làm tròn đến hàng phần chục nghìn là $0{,}0068$}{\True Trong ba phép đo trên, phép đo của bạn An chính xác nhất}
\loigiai{a) Đúng: $\delta_1\le\dfrac{0{,}5}{168}$.

b) Sai: $\delta_2\le\dfrac2{181}$ (sai số tương đối không âm).

c) Đúng: $\dfrac1{148}=0{,}006756\ldots\approx0{,}0068$.

d) Đúng: $\dfrac{0{,}5}{168}\approx0{,}0030<\dfrac1{148}\approx0{,}0068<\dfrac2{181}\approx0{,}0110$.}
\end{ex}
